\phantomsection\addcontentsline{toc}{section}{Circular Economy}
\ECchapter{20}{Circular Economy}{How can products remain useful for longer?}{ECRed}

\ECObjectives{
\begin{itemize}
\item Compare linear and circular product life cycles using delivered service as a reference.
\item Explain design for durability, repair, disassembly, refurbishment and remanufacturing.
\item Use life-cycle and multicriteria reasoning to avoid shifting environmental impacts.
\end{itemize}}

\begin{multicols}{2}
\begin{engineeringnews}
Manufacturers are developing repairable products, component take-back schemes and remanufacturing
lines. Circularity is not achieved by recycling alone: the highest value is often preserved when a
product, module or component remains in use with the least possible transformation.
\end{engineeringnews}

\begin{ecarticle}
The traditional linear model follows a simple sequence: extract resources, manufacture a product,
use it and discard it. This model loses materials, embodied energy and technical value at the end of
each product life. A circular economy aims to keep products and materials useful for longer through
maintenance, reuse, repair, refurbishment, remanufacturing and, finally, recycling.

Engineering decisions made during design strongly influence these options. A durable product needs
components that resist expected loads and environments. A repairable product needs accessible
fasteners, diagnostic information and available spare parts. Modular architecture allows a failed or
obsolete function to be replaced without discarding the complete system.

Disassembly must also be considered. Permanent adhesives and mixed materials may simplify production
but make separation difficult. Screws, clips and clearly identified polymers can improve recovery,
although they may add cost, mass or assembly time. Engineers therefore balance circularity with safety,
performance and industrial constraints.

Remanufacturing restores a used product to a defined level of performance. Components are inspected,
cleaned, repaired or replaced and then tested. This can preserve more value than material recycling,
but it requires traceability and a predictable return flow. Digital product passports can record
materials, maintenance history and previous use, helping operators choose the most suitable recovery
route.

Life-cycle assessment prevents misleading solutions. A heavier product may last longer but require
more material and transport energy. A replaceable battery can extend product life, yet the sealing
system must still meet safety requirements. Engineers compare resource use, energy, emissions and
service life across the complete system boundary. The objective is not merely to close a loop, but to
reduce total impact while still providing the required service.
\end{ecarticle}

\begin{tcolorbox}[title=\sffamily\bfseries Engineering Insight,colback=yellow!10,colframe=ECOrange]
Recycling is important, but it often destroys much of the manufacturing value already invested in a
component. Repair and reuse usually preserve more value when technical condition, safety and logistics
allow them. The shortest loop is not automatically best unless the recovered product still satisfies
its requirements.
\end{tcolorbox}

\begin{engineeringfocus}
\textbf{Circular design hierarchy}
\begin{center}
\resizebox{\linewidth}{!}{%
\begin{tikzpicture}[font=\scriptsize,>=Latex,node distance=3.2mm]
\node[draw=ECGreen,fill=ECGreen!10,rounded corners,align=center] (use) {use and\\maintain};
\node[draw=ECBlue,fill=ECLightBlue,rounded corners,align=center,right=of use] (rep) {repair and\\upgrade};
\node[draw=ECPurple,fill=ECPurple!10,rounded corners,align=center,right=of rep] (ref) {refurbish or\\remanufacture};
\node[draw=ECOrange,fill=ECOrange!10,rounded corners,align=center,below=of ref] (rec) {recycle\\materials};
\node[draw=ECRed,fill=ECRed!8,rounded corners,align=center,below=of rep] (loss) {energy recovery\\or disposal};
\draw[->,thick] (use)--(rep); \draw[->,thick] (rep)--(ref); \draw[->,thick] (ref)--(rec);
\draw[->,thick] (rec)--(loss); \draw[->,thick,dashed] (ref.north) to[bend left=28] (use.north);
\end{tikzpicture}%
}
\end{center}
Short loops usually preserve more product value and require less transformation.
\end{engineeringfocus}

\begin{keyfigures}
\begin{tabularx}{\linewidth}{@{}Xr@{}}
Highest-value strategy & continued use\\
Repair enabler & accessibility\\
Upgrade enabler & modularity\\
Recovery enabler & disassembly\\
Decision tool & life-cycle assessment
\end{tabularx}
\end{keyfigures}

\begin{tcolorbox}[title=\sffamily\bfseries Relative life-cycle impact,colback=white,colframe=ECRed]
\begin{center}
\begin{tikzpicture}
\begin{axis}[xbar,width=.94\linewidth,height=50mm,xlabel={Relative impact index},xmin=0,xmax=110,
symbolic y coords={Reuse,Remanufacture,Refurbishment,Repair,Linear use},
ytick=data,grid=major,ticklabel style={font=\scriptsize},nodes near coords]
\addplot table[x=impact,y=strategy,col sep=comma]{data/EC20/lifecycle.csv};
\end{axis}
\end{tikzpicture}
\end{center}
\footnotesize Illustrative comparison for the same delivered service; real results depend on
transport, energy mix, product lifetime, failure rate and replacement parts.
\end{tcolorbox}

\begin{vocabulary}
\begin{tabularx}{\linewidth}{@{}lX@{}}
linear economy & économie linéaire\\
durability & durabilité\\
repairability & réparabilité\\
disassembly & démontage\\
spare part & pièce détachée\\
refurbishment & reconditionnement\\
remanufacturing & remise à neuf industrielle\\
traceability & traçabilité\\
product passport & passeport produit\\
life-cycle assessment & analyse du cycle de vie
\end{tabularx}
\end{vocabulary}

\begin{applicationexercise}
A company compares two strategies for providing 1,000 hours of service. A linear product lasts 500 h
and has an impact index of 100 per product. A repairable version lasts 1,000 h and has an initial
impact index of 125; one repair adds an impact index of 15.
\begin{enumerate}
\item Determine how many linear products are required to provide 1,000 h of service.
\item Calculate the total impact index of the linear strategy.
\item Calculate the total impact index of the repairable strategy, including one repair.
\item Determine the percentage reduction achieved by the repairable strategy.
\item Calculate the break-even repair impact above which the circular strategy is no longer preferable.
\item Give two real-life parameters that should be added to the model.
\end{enumerate}
\end{applicationexercise}

\begin{questions}
\item Why does repair usually preserve more value than recycling?
\item Give three design choices that improve repairability.
\item What is the difference between refurbishment and remanufacturing?
\item Why can permanent adhesive create an end-of-life problem?
\item How does life-cycle assessment prevent impact shifting?
\item Why must alternatives be compared for the same delivered service?
\end{questions}

\begin{engineeringchallenge}
Redesign a cordless drill, laptop or coffee machine for a longer useful life. Define requirements for
lifetime, repair time, safety and cost. Propose a modular architecture, reversible joining methods,
replaceable components, diagnostics and a collection/remanufacturing process. Compare the original and
redesigned products with a weighted decision matrix and identify one rebound or impact-shifting risk.
\end{engineeringchallenge}

\begin{speaking}
Should manufacturers be legally required to provide spare parts and repair information? Discuss cost,
safety, innovation, ownership and environmental impact, then propose one measurable obligation.
\end{speaking}
\end{multicols}

\subsection*{Sources}
\ECsource{European Commission -- Circular Economy Action Plan}{https://environment.ec.europa.eu/strategy/circular-economy-action-plan_en}
\ECsource{Ellen MacArthur Foundation -- Circular Economy}{https://www.ellenmacarthurfoundation.org/topics/circular-economy-introduction/overview}